\documentclass[a4paper,12pt]{ctexart}
\usepackage{../teaching-style}

\begin{document}
\pagestyle{empty}

\maintitle{计算进阶：数的运算}
\begin{center}
  \large 适用于5\~{}初一 \quad 带讲+答案
\end{center}
\vspace{0.3cm}

\chaptertitle{第一讲\quad 有理数混合运算}

\sectiontitle{知识梳理}
\knowbox{
\textbf{1. 有理数的分类}\\
整数（正整数、0、负整数）和分数（正分数、负分数）统称有理数。\\
\textbf{2. 运算法则}\\
同号两数相加，取相同的符号，并把绝对值相加。\\
异号两数相加，取绝对值较大的符号，用较大的绝对值减去较小的绝对值。\\
减去一个数等于加上这个数的相反数。\\
\textbf{3. 运算顺序}\\
先乘方，再乘除，最后加减；有括号先算括号里的。}

\sectiontitle{典型例题}

\begin{example}
计算：$(-23)+17+(-15)+28$
\end{example}
\begin{solution}
原式 $=(-23-15)+(17+28)=(-38)+45=7$ \quad（同号结合）
\end{solution}

\begin{example}
计算：$(-81)\div\dfrac{9}{4}\times\dfrac{4}{9}\div(-16)$
\end{example}
\begin{solution}
原式 $=(-81)\times\dfrac{4}{9}\times\dfrac{4}{9}\times\left(-\dfrac{1}{16}\right)=81\times\dfrac{4}{9}\times\dfrac{4}{9}\times\dfrac{1}{16}=1$
\end{solution}

\begin{example}
计算：$-1^4-(1-0.5)\times\dfrac{1}{3}\times[2-(-3)^2]$
\end{example}
\begin{solution}
原式 $=-1-\dfrac{1}{2}\times\dfrac{1}{3}\times(2-9)=-1-\dfrac{1}{6}\times(-7)=-1+\dfrac{7}{6}=\dfrac{1}{6}$
\end{solution}

\sectiontitle{分层练习}

\begin{exercise}
\item 计算：
  \begin{subexercise}
    \item $(-15)+(-23)+34$
    \item $28-(-14)+(-18)$
    \item $(-12)-(-15)+(-8)-(-20)$
    \item $(-3.2)+4.5+(-1.8)+(-2.5)$
    \item $\left(-\dfrac{2}{3}\right)+\dfrac{1}{2}+\left(-\dfrac{1}{3}\right)+\dfrac{3}{2}$
    \item $(-2.5)+\left(-\dfrac{3}{4}\right)+\left(-\dfrac{1}{2}\right)+1.25$
  \end{subexercise}

\item 计算：
  \begin{subexercise}
    \item $(-8)\times(-7)\div(-2)$
    \item $(-36)\div(-9)\times(-3)$
    \item $(-24)\times\left(\dfrac{1}{2}-\dfrac{1}{3}+\dfrac{1}{4}\right)$
    \item $(-18)\div\left(\dfrac{1}{3}-\dfrac{1}{6}+\dfrac{1}{9}\right)$
    \item $(-5)\times(-4)+(-28)\div7$
    \item $(-3)^2\times(-2)+(-12)\div(-4)$
  \end{subexercise}

\item 计算：
  \begin{subexercise}
    \item $-2^2-(-3)^3\times(-1)^4$
    \item $(-1)^{10}+(-2)^3\div(-4)$
    \item $(-2)^2-(-5^2)\times(-1)^5$
    \item $-1^4-\dfrac{1}{6}\times[2-(-3)^2]$
    \item $(-3)^2-(-3)\times\dfrac{1}{3}\times\left[1-(-2)^3\right]$
    \item $(-1)^5\times\left[4\dfrac{2}{3}\div(-4)+(-1\dfrac{1}{4})\times(-0.4)\right]\div\left(-\dfrac{1}{3}\right)$
  \end{subexercise}
\end{exercise}

\newpage

\chaptertitle{第二讲\quad 绝对值与数轴}

\sectiontitle{知识梳理}
\knowbox{
\textbf{1. 绝对值的定义}\\
数轴上表示数$a$的点到原点的距离叫做$a$的绝对值，记作$|a|$。\\
$|a|=\begin{cases}
a & (a>0)\\
0 & (a=0)\\
-a & (a<0)
\end{cases}$\\
\textbf{2. 绝对值的性质}\\
$|a|\ge0$；$|ab|=|a||b|$；$\left|\dfrac{a}{b}\right|=\dfrac{|a|}{|b|}(b\neq0)$。\\
\textbf{3. 数轴}\\
规定了原点、正方向和单位长度的直线。数轴上右边的数总比左边的大。}

\sectiontitle{典型例题}

\begin{example}
已知$|x|=3$，$|y|=5$，且$|x-y|=y-x$，求$x+y$的值。
\end{example}
\begin{solution}
由$|x|=3$得$x=\pm3$，$|y|=5$得$y=\pm5$。\\
$|x-y|=y-x$说明$y-x\ge0$即$y\ge x$。\\
当$y=5$时，$x=\pm3$均满足，$x+y=8$或$2$。\\
当$y=-5$时，$y< x$，不满足。\\
故$x+y=8$或$2$。
\end{solution}

\begin{example}
化简：$|x-1|+|x-3|$（分别讨论$x<1$，$1\le x<3$，$x\ge3$三种情况）
\end{example}
\begin{solution}
当$x<1$时，原式$=(1-x)+(3-x)=4-2x$\\
当$1\le x<3$时，原式$=(x-1)+(3-x)=2$\\
当$x\ge3$时，原式$=(x-1)+(x-3)=2x-4$
\end{solution}

\sectiontitle{分层练习}

\begin{exercise}
\item 计算：
  \begin{subexercise}
    \item $|{-5}|+|{-3}|$
    \item $|{-8}|-|{-2}|$
    \item $|{-3}|\times|{-4}|$
    \item $|{-12}|\div|{-3}|$
    \item $|-2^2|+|-3^2|$
    \item $|-5|-|{-3}^2|$
  \end{subexercise}

\item 比较大小：
  \begin{subexercise}
    \item $-\dfrac{2}{3}$ 和 $-\dfrac{3}{4}$
    \item $-0.618$ 和 $-\dfrac{5}{8}$
    \item $-|{-3}|$ 和 $-(-3)$
    \item $-\dfrac{5}{7}$ 和 $-\dfrac{3}{5}$
  \end{subexercise}

\item 已知$|x|=2$，$|y|=3$，求$x+y$的值。

\item 已知$|x-2|+|y+3|=0$，求$x$、$y$的值。

\item 化简（分段讨论）：
  \begin{subexercise}
    \item $|x+1|$
    \item $|x-2|+|x+3|$
    \item $|2x-1|$
    \item $|x+1|-|x-2|$
  \end{subexercise}

\item 若$a$、$b$互为相反数，$c$、$d$互为倒数，$|m|=2$，求$\dfrac{a+b}{m}+m^2-cd$的值。
\end{exercise}

\newpage

\chaptertitle{第三讲\quad 平方根与立方根}

\sectiontitle{知识梳理}
\knowbox{
\textbf{1. 平方根}\\
若$x^2=a(a\ge0)$，则$x$叫做$a$的平方根，记作$x=\pm\sqrt{a}$。\\
正数有两个平方根，互为相反数；0的平方根是0；负数没有平方根。\\
\textbf{2. 算术平方根}\\
正数$a$的正的平方根叫做算术平方根，记作$\sqrt{a}$。\\
\textbf{3. 立方根}\\
若$x^3=a$，则$x$叫做$a$的立方根，记作$x=\sqrt[3]{a}$。\\
任何实数都有唯一确定的立方根。}

\sectiontitle{典型例题}

\begin{example}
求下列各数的平方根和算术平方根：\\
(1) $144$ \qquad (2) $\dfrac{25}{49}$ \qquad (3) $0.81$ \qquad (4) $(-4)^2$
\end{example}
\begin{solution}
(1) $\pm\sqrt{144}=\pm12$，算术平方根为$12$。\\
(2) $\pm\sqrt{\dfrac{25}{49}}=\pm\dfrac{5}{7}$，算术平方根为$\dfrac{5}{7}$。\\
(3) $\pm\sqrt{0.81}=\pm0.9$，算术平方根为$0.9$。\\
(4) $\pm\sqrt{(-4)^2}=\pm\sqrt{16}=\pm4$，算术平方根为$4$。
\end{solution}

\begin{example}
计算：$\sqrt{25}-\sqrt[3]{27}+\sqrt{\dfrac{1}{4}}-\sqrt[3]{-8}$
\end{example}
\begin{solution}
原式 $=5-3+\dfrac{1}{2}-(-2)=5-3+0.5+2=4.5$
\end{solution}

\sectiontitle{分层练习}

\begin{exercise}
\item 求下列各数的平方根：
  \begin{subexercise}
    \item $169$
    \item $0.25$
    \item $\dfrac{64}{81}$
    \item $2\dfrac{1}{4}$
    \item $(-6)^2$
    \item $10^4$
  \end{subexercise}

\item 求下列各数的立方根：
  \begin{subexercise}
    \item $27$
    \item $-64$
    \item $0.125$
    \item $-\dfrac{8}{27}$
    \item $216$
    \item $-1$
  \end{subexercise}

\item 计算：
  \begin{subexercise}
    \item $\sqrt{81}+\sqrt{25}$
    \item $\sqrt{144}-\sqrt{49}$
    \item $\sqrt{0.09}+\sqrt{0.16}$
    \item $\sqrt[3]{8}+\sqrt[3]{-27}$
    \item $\sqrt{36}-\sqrt[3]{-64}+\sqrt{\dfrac{1}{4}}$
    \item $\sqrt[3]{125}-\sqrt{100}+\sqrt[3]{0.001}$
  \end{subexercise}

\item 求下列各式中$x$的值：
  \begin{subexercise}
    \item $x^2=49$
    \item $4x^2=25$
    \item $(x-1)^2=9$
    \item $x^3=-27$
    \item $8x^3+27=0$
    \item $(x+2)^3=64$
  \end{subexercise}
\end{exercise}

\newpage

\chaptertitle{第四讲\quad 实数混合运算}

\sectiontitle{知识梳理}
\knowbox{
\textbf{1. 实数的分类}\\
有理数（有限小数或无限循环小数）和无理数（无限不循环小数）。\\
常见无理数：$\pi$，$e$，$\sqrt{2}$，$\sqrt{3}$等。\\
\textbf{2. 实数的运算}\\
有理数的运算法则和运算律对实数同样适用。\\
$\sqrt{a}\cdot\sqrt{b}=\sqrt{ab}\;(a\ge0,b\ge0)$，$\dfrac{\sqrt{a}}{\sqrt{b}}=\sqrt{\dfrac{a}{b}}\;(a\ge0,b>0)$。\\
$\sqrt{a^2}=|a|$；$(\sqrt{a})^2=a\;(a\ge0)$。}

\sectiontitle{典型例题}

\begin{example}
计算：$\sqrt{12}-3\sqrt{\dfrac{1}{3}}+2\sqrt{48}$
\end{example}
\begin{solution}
原式 $=2\sqrt{3}-3\times\dfrac{\sqrt{3}}{3}+2\times4\sqrt{3}=2\sqrt{3}-\sqrt{3}+8\sqrt{3}=9\sqrt{3}$
\end{solution}

\begin{example}
计算：$(\sqrt{3}+2)(\sqrt{3}-2)+(\sqrt{5}+1)^2$
\end{example}
\begin{solution}
原式 $=(3-4)+(5+2\sqrt{5}+1)=-1+6+2\sqrt{5}=5+2\sqrt{5}$
\end{solution}

\sectiontitle{分层练习}

\begin{exercise}
\item 计算：
  \begin{subexercise}
    \item $\sqrt{8}+\sqrt{18}-\sqrt{32}$
    \item $\sqrt{27}-\sqrt{\dfrac{1}{3}}+\sqrt{12}$
    \item $\sqrt{50}-3\sqrt{2}+2\sqrt{8}$
    \item $\sqrt{75}+2\sqrt{\dfrac{1}{3}}-3\sqrt{48}$
    \item $2\sqrt{20}-3\sqrt{45}+\sqrt{80}$
    \item $\sqrt{54}-\sqrt{24}+2\sqrt{\dfrac{2}{3}}$
  \end{subexercise}

\item 计算：
  \begin{subexercise}
    \item $\sqrt{3}\times\sqrt{6}$
    \item $\sqrt{2}\times\sqrt{8}$
    \item $\sqrt{27}\div\sqrt{3}$
    \item $\sqrt{50}\div\sqrt{2}$
    \item $\sqrt{12}\times\sqrt{3}-5$
    \item $(\sqrt{5}+1)(\sqrt{5}-1)$
  \end{subexercise}

\item 计算：
  \begin{subexercise}
    \item $(\sqrt{2}+1)^2$
    \item $(\sqrt{3}-2)^2$
    \item $(\sqrt{7}+\sqrt{5})(\sqrt{7}-\sqrt{5})$
    \item $(\sqrt{2}+\sqrt{3})^2$
    \item $(\sqrt{5}-\sqrt{3})^2$
    \item $(2\sqrt{3}+3\sqrt{2})(2\sqrt{3}-3\sqrt{2})$
  \end{subexercise}
\end{exercise}

\newpage

\chaptertitle{综合闯关}

\begin{exercise}
\item 计算：$(-12)+(-23)+34+(-18)$
\item 计算：$(-2)^3\times(-3)+(-36)\div(-4)$
\item 计算：$(-1)^4-(-5\dfrac{1}{2})\times\dfrac{4}{11}+(-2)^3\div|-4|$
\item 已知$|x-3|+|y+2|=0$，求$x+y$的值。
\item 若$a$、$b$互为相反数，$c$、$d$互为倒数，$|m|=3$，求$\dfrac{a+b}{m}+m^2-cd$的值。
\item 求$\sqrt{169}-\sqrt[3]{125}+\sqrt{0.81}-\sqrt[3]{-0.008}$
\item 求$x$：$4(x-1)^2-9=0$
\item 计算：$\sqrt{48}-2\sqrt{3}+3\sqrt{\dfrac{1}{3}}$
\item 计算：$(\sqrt{5}+2)(\sqrt{5}-2)-(\sqrt{3}-1)^2$
\item 比较大小：$-\dfrac{5}{6}$ 和 $-\dfrac{6}{7}$
\end{exercise}

\newpage

\chaptertitle{参考答案}

\textbf{第一讲}
\noindent\\
1.(a) $-4$ \quad (b) $24$ \quad (c) $15$ \quad (d) $-3$ \quad (e) $1$ \quad (f) $-2.5$\\
2.(a) $-28$ \quad (b) $-12$ \quad (c) $-10$ \quad (d) $-324$ \quad (e) $16$ \quad (f) $21$\\
3.(a) $23$ \quad (b) $3$ \quad (c) $-21$ \quad (d) $\dfrac{7}{6}$ \quad (e) $30$ \quad (f) $\dfrac{11}{4}$\\

\textbf{第二讲}
\noindent\\
1.(a) $8$ \quad (b) $6$ \quad (c) $12$ \quad (d) $4$ \quad (e) $13$ \quad (f) $-4$\\
2.(a) $-\dfrac{2}{3}>-\dfrac{3}{4}$ \quad (b) $-0.618>-\dfrac{5}{8}$ \quad (c) $-|{-3}|<-(-3)$ \quad (d) $-\dfrac{5}{7}<-\dfrac{3}{5}$\\
3. $\pm5$，$\pm1$ \quad 4. $x=2$，$y=-3$\\
5.(a) $\begin{cases}x+1&(x\ge-1)\\-x-1&(x<-1)\end{cases}$ \quad (b) $\begin{cases}-2x-1&(x<-3)\\5&(-3\le x<2)\\2x+1&(x\ge2)\end{cases}$\\
6. $3$\\

\textbf{第三讲}
\noindent\\
1.(a) $\pm13$ \quad (b) $\pm0.5$ \quad (c) $\pm\dfrac89$ \quad (d) $\pm\dfrac32$ \quad (e) $\pm6$ \quad (f) $\pm100$\\
2.(a) $3$ \quad (b) $-4$ \quad (c) $0.5$ \quad (d) $-\dfrac23$ \quad (e) $6$ \quad (f) $-1$\\
3.(a) $14$ \quad (b) $5$ \quad (c) $0.7$ \quad (d) $-1$ \quad (e) $\dfrac{25}{2}$ \quad (f) $2.1$\\
4.(a) $\pm7$ \quad (b) $\pm\dfrac52$ \quad (c) $4$或$-2$ \quad (d) $-3$ \quad (e) $-\dfrac32$ \quad (f) $2$\\

\textbf{第四讲}
\noindent\\
1.(a) $\sqrt2$ \quad (b) $\dfrac{8\sqrt3}{3}$ \quad (c) $4\sqrt2$ \quad (d) $-5\sqrt3$ \quad (e) $-2\sqrt5$ \quad (f) $\dfrac{8\sqrt6}{3}$\\
2.(a) $3\sqrt2$ \quad (b) $4$ \quad (c) $3$ \quad (d) $5$ \quad (e) $1$ \quad (f) $4$\\
3.(a) $3+2\sqrt2$ \quad (b) $7-4\sqrt3$ \quad (c) $2$ \quad (d) $5+2\sqrt6$ \quad (e) $8-2\sqrt{15}$ \quad (f) $-6$\\

\textbf{综合闯关}
\noindent\\
1. $-19$ \quad 2. $33$ \quad 3. $4$ \quad 4. $1$ \quad 5. $8$ \quad 6. $3.9$ \quad 7. $x=\dfrac52$或$-\dfrac12$ \quad 8. $\dfrac{11\sqrt3}{3}$ \quad 9. $2\sqrt3-3$ \quad 10. $-\dfrac56>-\dfrac67$

\end{document}
